\chapter{绪论}
\label{chp:instruction}

\section{研究背景及意义}
\label{sec:background}

随着人工智能与机器人技术的蓬勃发展，具身智能作为连接数字世界与物理世界的桥梁，正成为国家战略布局的核心领域。工业和信息化部在二零二三年印发的《人形机器人创新发展指导意见》中明确指出，需推动人形机器人产业高质量发展，培育形成新质生产力，并设定了到二零二五年初步建立创新体系、突破“大脑、小脑、肢体”等关键技术的发展目标\cite{miit2023humanoid}。同时，国际机器人联合会（IFR）发布的《World Robotics 2024》报告显示，二零二三年全球工业机器人安装量连续第三年超过五十万台，其中中国工厂中运转的机器人数量达创纪录的一百七十万台\cite{ifr2024world}。在这一宏观背景下，机械臂作为智能制造和家庭服务等领域的核心执行部件，其自主操作能力被视为实现具身智能落地的关键一环\cite{shen2025embodied}。赋予机械臂在非结构化环境中理解自然语言指令并执行复杂操作的能力，不仅是学术界的前沿挑战，更是推动智能机器人向通用化、智能化迈进的必由之路\cite{lan2024mobile}。

在机械臂自主操作任务中，环境感知是实现智能决策与精准执行的先决条件。然而，真实世界中的非结构化环境往往伴随着光照变化、物体层叠遮挡以及视角受限等复杂因素，给机器人的视觉感知带来了巨大挑战\cite{wang2025embodied}。传统的视觉感知方法多依赖于人工设计的特征提取器或特定场景下的标定模型，难以在动态变化的环境中保持鲁棒性。为克服这些限制，基于深度学习的视觉感知与抓取算法逐渐成为研究热点，通过端到端的数据驱动方式显著提升了机械臂在特定任务中的表现\cite{mahler2017dexnet, fang2020graspnet}。然而，这类方法往往依赖大规模标注数据集进行离线训练，面对开放场景中物体的长尾分布时泛化性能欠佳，难以满足真实环境对多任务、多目标操作的复杂需求。

在控制与决策层面，尽管机械臂自主操作技术展现出巨大的应用潜力，但传统控制方法在应对复杂开放场景时仍面临显著局限。以快速扩展随机树（RRT*）为代表的传统运动规划算法高度依赖预定义的精确环境模型，难以适应现实场景中物体的位姿变化以及未知障碍物等挑战\cite{lavalle1998rapidly}。为了实现更灵活的控制，深度强化学习（DRL）被广泛引入机械臂操作领域，为自适应决策带来了新契机\cite{lillicrap2015continuous, haarnoja2018sac}。然而，强化学习在处理高维连续动作空间时面临严重的样本效率低下问题，且在长视野任务中往往受困于稀疏奖励的挑战，难以在真实物理环境中高效收敛。另一方面，基于行为克隆（Behavior Cloning, BC）的模仿学习方法虽然能够直接从专家演示中学习策略，但其在面对分布外状态时容易产生协变量偏移（Covariate Shift）和复合误差，导致策略在执行长序列任务时迅速崩溃\cite{ross2011dagger}。此外，传统的行为克隆方法在将多模态特征映射为动作时，往往假设动作分布是单峰的，这与真实世界中同一任务存在多种可行操作路径的多模态特性相悖，极大限制了策略的鲁棒性和适应性。

近年来，去噪扩散概率模型（DDPM）在生成领域的成功为解决上述模仿学习难题提供了全新思路\cite{ho2020denoising}。将扩散模型引入机器人操作领域作为动作生成专家（如Diffusion Policy），能够有效建模复杂的多模态动作分布，生成鲁棒且适应性强的动作序列\cite{chi2023diffusion}。扩散策略通过迭代去噪过程，不仅能够平滑地处理高维连续动作空间，还能在一定程度上缓解行为克隆中的复合误差问题。与此同时，大语言模型（LLM）和视觉语言模型（VLM）的突破性进展为机械臂操作带来了新的认知范式。研究人员开始探索将视觉感知与语言理解相融合，构建视觉-语言-动作（VLA）模型，以实现基于自然语言指令的机器人控制。例如，SayCan通过语言模型评估机器人技能的可行性\cite{ahn2022saycan}；RT-1和RT-2将视觉语言模型直接用于端到端机器人控制\cite{brohan2022rt1, brohan2023rt2}；OpenVLA和RoboFlamingo则基于开源大模型构建了高效的机器人操作框架\cite{kim2024openvla, li2024roboflamingo}。然而，现有的端到端VLA模型多采用统一的架构处理多模态输入，这种高度耦合的设计虽然简化了系统结构，但也带来了特征表征不充分、模型可解释性差以及训练成本高昂等问题。特别是在面对复杂的精细操作任务时，端到端模型难以显式地建立动力学约束，容易导致机械臂动作出现关节超限或轨迹震荡。

针对上述挑战，本文提出了一种基于解耦视觉-语言特征与扩散策略的机械臂模仿学习操作算法。该算法创新性地采用“特征解耦-扩散生成”的双层架构，突破了传统端到端模型的不可解释性瓶颈。在特征理解阶段，本研究摒弃了高度耦合的视觉语言模型，转而采用独立的视觉编码器（如CLIP或ViT）提取高维空间特征\cite{radford2021clip, dosovitskiy2020vit}，并结合纯语言大模型解析自然语言指令，从而实现视觉感知与语义理解的深度解耦与精准对齐。在动作生成阶段，本文引入扩散模型作为模仿学习的动作专家，直接从关节空间生成动作序列，规避了逆运动学求解的多解性问题。这种基于解耦特征驱动扩散策略的方法，不仅充分发挥了扩散模型在多模态动作分布建模上的优势，还显著提升了机械臂在非结构化环境下面对多任务指令的泛化能力与操作精度，为破解具身智能领域中语义-动作联合推理难题提供了新的理论支撑与实践路径。

% 占位符：后续在此处添加系统架构图
% \begin{figure}[htbp]
%   \centering
%   \includegraphics[width=0.8\linewidth]{figures/content/placeholder_architecture}
%   \caption{基于解耦视觉-语言特征与扩散策略的机械臂操作算法架构图}
%   \label{fig:architecture}
% \end{figure}
